\scanpage{145}
\mspage{90}
% archive sheet number: 144

\unsure{.} On se place donc dans l'hyp. quillénienne sur les $i_{\deleted{\unsure{Y}}X}$. On choisit $C(X),i_X$ pour chaque $X$ \inserted{dans $\mathcal M_{/S}$}, et on pose comme tantôt

\[
\gamma(X)=C(X)\qquad\text{comme objet de }(\mathcal M_{/S}\unsure{,|)}W_S^{-1}.
\]

Montrons que c'est fonctoriel en $X$. Soit

\[
u:X\unsure{-|\to}X'
\]

dans $\mathcal M_{/S}$, considérons le diagramme anodin \inserted{(\unsure{comm.})}

\diagram{\begin{tikzcd}[ampersand replacement=\&,column sep=1.4em,row sep=1.3em]
\deleted{\unsure{Y}}X\arrow[rr,"u"]\arrow[d,"i_X"']\arrow[drr,dashed,"{v=i_{X'}u}"]\&\&X'\arrow[d,"i_{X'}"]\\
C(X)\arrow[rr,dotted,"\bar u"]\arrow[dr]\&\& C(X')\arrow[dl,"{\in F}"]\\
\&S\&,
\end{tikzcd}}

\margin{Lemme: soit
\diagram{\begin{tikzcd}[ampersand replacement=\&,column sep=0.6em,row sep=0.55em,font=\scriptsize]
X\arrow[r,-{Bar},"u"]\arrow[d,dash,"i"']\&X'\arrow[d]\\
C\arrow[r,dotted]\&C'
\end{tikzcd}}
dans $\mathcal M_{/S}$ avec $i\in\mathrm{Cof}^{\deleted{\unsure{a}}}_0$ et $(C'\unsure{-|\to}S)\in\unsure{F|\mathcal F}\deleted{\unsure{=\mathrm{Fib}}}$ \inserted{$(=\mathrm{Fib})$}, alors la flèche en pointillés existe, \inserted{dans $\mathcal M_{/S}$}, de façon à rendre le carré comm.}

complétons le par la flèche en \inserted{gros} pointillés, $v=i_{X'}u$ (rendant le triangle sup. commutatif), l'hypothèse sur $i_X$ implique que la flèche $\bar u$ en pointillés fins existe, de façon à avoir encore deux triangles commutatifs. Mais ça signifie exactement qu'il y a un \uline{$S$-morphisme}, $\bar u$, i.e. une flèche dans $\deleted{\unsure{\mathrm C}}\mathcal M_{/S}$, rendant commutatif le carré

\diagram{\begin{tikzcd}[ampersand replacement=\&]
X\arrow[r,"u"]\arrow[d,dash,"i_X"']\& X'\arrow[d,"i_{X'}"]\\
C(X)\arrow[r,"\bar u"]\&C(X'),
\end{tikzcd}}
